\section{Intersection amplitudes and the comparison model}\label{exact-amplitudes-and-the-harmonic-model}

We deform the Weil--Petersson intersection brackets into pure $\psi$
intersection numbers and derive a positive recursion along the deformation.
We call the resulting functions shifted amplitudes. The harmonic model
introduced below is an explicit, solvable comparison for this recursion.
Section~\ref{uniform-comparison} compares the full recursion with that
model, uniformly in the degrees and in the deformation parameter.
This comparison is the common starting point for the volume and
coefficient asymptotics.

\subsection{Intersection identities and normalization}\label{exact-amplitudes-standard-intersection-identities-and-normalization}

We use the following notation for a stable pair \((g,n)\) and a degree
tuple \(\mathbf d\): \(h\) is the complexity, \(N\) the dimension of
\(\overline{\mathcal M}_{g,n}\), and the deficit \(m\) is the number of
\(\kappa_1\) factors needed to fill up the dimension. Write \[
h=2g-2+n,\qquad N=3g-3+n,\qquad
D=\sum_{i=1}^n d_i,\qquad m=N-D.
\] Throughout the comparison theorem \(n\ge1\), the pair is stable, and
\(0\le D\le N\). A tuple with a negative degree or an unstable pair has
amplitude zero.

Recall \(X\), \(a=X'>0\) on \([0,y_0)\) and the kernels \(K_Y\) and
\(H_Y\) of \eqref{eq:P1}, and write \(r_Y=r(Y)\) for the first positive
zero of \(K_Y\). The properties of these kernels used below were
established in Section~\ref{the-positive-compact-family-of-kernels}. Put \begin{equation}
\bar H_Y(z)=\frac{H_Y(r_Yz)}a,\qquad
\bar K_Y(z)=\frac{K_Y(r_Yz)}a
      =\bar H_Y(z)+\frac23z\bar H_Y'(z),\qquad
\theta_\ell(Y)=[z^\ell]\bar K_Y(z)^{-1}.
\label{eq:H2}
\end{equation} At \(Y=y_0\) these have the continuous definitions \begin{equation}
\bar H_{y_0}(z)=1-\frac35z,\qquad \bar K_{y_0}(z)=1-z,\qquad
\theta_\ell(y_0)=1.
\label{eq:H3}
\end{equation} In particular \(\theta_0=1\). These are the rescaled kernels and
coefficients of \eqref{eq:P9} and \eqref{eq:P11}, recalled here together
with their endpoint values.

We use the following facts from intersection theory. The ordinary string,
dilaton, and Dijkgraaf--Verlinde--Verlinde (DVV) identities hold for Witten's psi
intersections~\cite{Witten,Kontsevich,DVV}, with
\(\langle\tau_0^3\rangle_0=1\) and \(\langle\tau_1\rangle_1=1/24\). The normalized forms used here are written in Appendix~\ref{app:primitive-bounds}. For
pure $\psi$ intersection numbers we use the normalization \begin{equation}
C_g(\mathbf e)=
\frac{4^g\prod_i(2e_i+1)!!}
 {3^{\,2g-2+q}\Gamma(2g-2+q)}
\left\langle\prod_{i=1}^q\tau_{e_i}\right\rangle_g,
\qquad \sum_i e_i=3g-3+q.
\label{eq:H4}
\end{equation} For positive degrees we use only
\[
 \frac16\le C_g(\mathbf e)\le1.
\]
Appendix~\ref{app:primitive-bounds} proves these coarse bounds directly
from DVV and the published two-point estimate. The argument adapts the
recursion estimates of Guo, Yang and Zagier~\cite{GYZ}; their stronger
lower bound $C_g(3g-2)\le C_g(\mathbf e)$ is not needed here.

With \(m=3g-3+n-|\mathbf d|\), our mixed brackets are
\[
 [\tau_{d_1}\cdots\tau_{d_n}]_{g,n}
 =\prod_i(2d_i+1)!!\,4^{|\mathbf d|}\frac{(2\pi^2)^m}{m!}
 \int_{\overline{\mathcal M}_{g,n}}\kappa_1^m\prod_i\psi_i^{d_i}.
\]
Here \(\kappa_1\) is the logarithmic Mumford class, \(\psi_i\) is the cotangent class, and \([\omega_{\mathrm{WP}}]=2\pi^2\kappa_1\)~\cite{Wolpert}; see also~\cite[Section~2]{MZ}. Thus
\[
 V_{g,n}(2\mathbf L)=\sum_{|\mathbf d|\le3g-3+n}
 [\tau_{d_1}\cdots\tau_{d_n}]_{g,n}
 \prod_i\frac{L_i^{2d_i}}{(2d_i+1)!}.
\]
Angle brackets below denote unweighted pure-$\psi$ integrals. Intersections of the wrong dimension and negative-degree amplitudes are zero. Unstable factors in the positive recursion are zero; explicitly specified genus-zero initial derivatives are kept separate.

The mixed-kappa/psi transform~\cite{KMZ,LX} gives, at
\(Y=0\), \begin{equation}
\frac{[\tau_{d_1}\cdots\tau_{d_n}]_{g,n}}
 {4^D(2\pi^2)^m}
=\prod_i(2d_i+1)!!
\sum_{\substack{\ell_k\ge0,\ k\ge2\\
                 \sum_{k\ge2}(k-1)\ell_k=m}}
\left\langle\tau_{\mathbf d}\prod_{k\ge2}\tau_k^{\ell_k}\right\rangle_g
\prod_{k\ge2}\frac{((-1)^k/(k-1)!)^{\ell_k}}{\ell_k!}.
\label{eq:H5}
\end{equation}
Since \(X'(0)=1\) and \(-X^{(k)}(0)=(-1)^k/(k-1)!\), the weights in
\eqref{eq:H5} are the derivatives \(-X^{(k)}(0)\), \(k\ge2\). The shifted
amplitudes replace them by the derivatives \(-X^{(k)}(Y)\) at a point
\(Y\). For general \(Y<y_0\), define the shifted amplitudes by the finite sum
\begin{equation}
U^Y_{g,n}(\mathbf d)=\prod_i(2d_i+1)!!
\sum_{\substack{\ell_k\ge0,\ k\ge2\\
                 \sum(k-1)\ell_k=m}}
\left\langle\tau_{\mathbf d}\prod_{k\ge2}\tau_k^{\ell_k}\right\rangle_g
\frac{\prod_{k\ge2}(-X^{(k)}(Y))^{\ell_k}/\ell_k!}
 {a^{h+\sum\ell_k}}.
\label{eq:H6}
\end{equation} The denominator is obtained by summing degree-one background
insertions with time \(t_1=1-a\): the dilaton equation gives the factor
\(\sum_{j\ge0}(h+\sum\ell_k)_j t_1^j/j!=a^{-h-\sum\ell_k}\). Define the
normalized amplitudes \begin{equation}
\mathcal U_{g,n}(\mathbf d;Y)=a^h r_Y^m U^Y_{g,n}(\mathbf d).
\label{eq:H7}
\end{equation} They extend to \(Y=y_0\) through their positive recursion below.

When all degrees vanish, the shifted amplitudes are derivatives of the
cusp generating function. This is the familiar shift of the Kontsevich--Witten (KW) times
along the string flow~\cite{IZ,KMZ}; the identification of $\Phi_g$ with the
KW free energy at the shifted times $t_0=x$, $t_1=0$, $t_k=(-1)^k/(k-1)!$
is recorded in~\cite[Remark~1]{Zog98}, where it is attributed to Dubrovin, and the corresponding deformation in
$(t_0,t_1)$ is used in~\cite[Section~2.4]{OS}. We recall it in the present
normalization. Start with KW
times \(t_0=x,\ t_1=0,\ t_k=-X^{(k)}(0)\), \(k\ge2\). The finite formal
string flow sends them to \[
t_0'=x-X(Y),\qquad t_1'=1-X'(Y),\qquad
t_k'=-X^{(k)}(Y)\quad(k\ge2).
\] Its vector field is
\(\sum_{k\ge0}t_{k+1}'\partial_{t_k'}-\partial_{t_0'}\). The string
equation makes the genus-\(g\) free energy invariant for \(g\ge1\); its
genus-zero anomaly disappears after three \(t_0\) derivatives. Since
\(\partial_{t_0}\) commutes with this flow, setting \(x=X(Y)\) gives \begin{equation}
U^Y_{g,n}(0^n)=\Phi_g^{(n)}(x),\qquad
\Phi_g(x)=\sum_{n\ge0}\frac{V_{g,n}}{(2\pi^2)^{3g-3+n}}\frac{x^n}{n!},
\label{eq:H8}
\end{equation} whenever the amplitude is stable. Unmarked \(\Phi_g\) is used only
for \(g\ge2\). The unstable genus-zero derivatives are separate initial
terms and are not substituted into the positive recursion.

\subsection{The positive recursion}\label{deriving-the-exact-positive-recursion}

We single out the first marking and call it the \emph{pivot}. Write
\(\langle\!\langle\tau_{d_1}\cdots\tau_{d_q}\rangle\!\rangle_g
=\prod_i(2d_i+1)!!\,\langle\tau_{d_1}\cdots\tau_{d_q}\rangle_g\).
In this normalization the Dijkgraaf--Verlinde--Verlinde equation for the
pivot reads, with \((-1)!!=1\),
\begin{equation*}
\begin{split}
\langle\!\langle\tau_{d_1}\tau_{d_2}\cdots\tau_{d_q}\rangle\!\rangle_g
={}&\sum_{j=2}^q(2d_j+1)\,
 \langle\!\langle\tau_{d_1+d_j-1}\textstyle\prod_{i\ne1,j}\tau_{d_i}\rangle\!\rangle_g
 +\frac12\sum_{u+v=d_1-2}
 \langle\!\langle\tau_u\tau_v\tau_{d_2}\cdots\tau_{d_q}\rangle\!\rangle_{g-1}\\
&+\frac12\sum_{u+v=d_1-2}\ \sum_{\substack{g_1+g_2=g\\ I\sqcup J=\{2,\ldots,q\}}}
 \langle\!\langle\tau_u\tau_I\rangle\!\rangle_{g_1}
 \langle\!\langle\tau_v\tau_J\rangle\!\rangle_{g_2}.
\end{split}
\end{equation*}
For \(d_1=0\) and \(d_1=1\) this is the string and the dilaton equation.
In the amplitudes \eqref{eq:H6}, besides the \(n\) external markings, the
intersection numbers contain \emph{background} insertions: \(\tau_k\) with
weight \(t_k=-X^{(k)}(Y)\) for \(k\ge2\), and, after the dilaton sum that
produced the factor \(a^{-h-\sum\ell_k}\), also \(\tau_1\) with weight
\(t_1=1-a\). Apply the displayed equation with the pivot \(\tau_{d_1}\)
and sum over the background. Three kinds of terms appear.

\begin{enumerate}[label=(\alph*),leftmargin=*]
\item \emph{The pivot merges with an external marking $j\ge2$.} By the
first term above, this gives \((2d_j+1)\,U^Y_{g,n-1}(d_1+d_j-1,\ldots)\),
with the background unchanged.
\item \emph{The pivot merges with a background insertion $\tau_k$,
$k\ge1$.} Background insertions carry no double factorial, so the
factor \(2k+1\) of the first term is divided by \((2k+1)!!\) and becomes
\(1/(2k-1)!!\). Choosing which of the
\(\ell_k\) copies is used turns \(t_k^{\ell_k}/\ell_k!\) into
\(t_k\cdot t_k^{\ell_k-1}/(\ell_k-1)!\). The result is
\[
\frac{t_k}{(2k-1)!!}\,U^Y_{g,n}(d_1+k-1,d_2,\ldots,d_n).
\]
\item \emph{Contractions.} The two remaining terms give the nonseparating
and separating contributions with coefficient \(1/2\). In the separating
case the background splits between the two components with the binomial
factors of exponential generating functions, so the term is a product of
two amplitudes.
\end{enumerate}

Let \(E\) be the operator raising the pivot degree by one. The terms of
type~(b) form \((1-a)U^Y+\sum_{k\ge2}t_kE^{k-1}U^Y/(2k-1)!!\). Moving them
to the left-hand side and recalling \(t_k=-X^{(k)}(Y)\) gives
\[
\left(a+\sum_{k\ge2}\frac{X^{(k)}(Y)}{(2k-1)!!}E^{k-1}\right)U^Y
=K_Y(E)\,U^Y=\bigl(\text{terms of types (a) and (c)}\bigr),
\]
where we compared with the series~\eqref{eq:P1} for \(K_Y\). For a fixed
parent, raising the pivot degree more than \(m\) times leaves the
dimension range, so \(E^{m+1}U^Y=0\) and \(K_Y(E)\) has the finite inverse
\(\sum_{\ell\le m}[z^\ell]K_Y(z)^{-1}E^\ell\). This is how the Taylor
coefficients of \(1/K_Y\) enter. Components of the wrong dimension vanish,
and for external complexity \(h\ge2\) no exceptional initial tuple occurs.

After \eqref{eq:H7}, the exact recursion is \begin{equation}
\begin{aligned}
\mathcal U_{g,n}(\mathbf d)={}&
\sum_{j=2}^n(2d_j+1)\sum_{\ell=0}^m\theta_\ell
 \mathcal U_{g,n-1}(d_1+d_j+\ell-1,\mathbf d_{\widehat{1,j}})\\
&+\frac12\sum_{\ell=0}^m\theta_\ell
 \sum_{u+v=d_1+\ell-2}
 \mathcal U_{g-1,n+1}(u,v,\mathbf d_{\widehat1})\\
&+\frac12\sum_{\ell=0}^m\theta_\ell
 \sum_{u+v=d_1+\ell-2}
 \sum_{\substack{g_1+g_2=g\\I\sqcup J=\{2,\ldots,n\}}}
 \mathcal U_{g_1,|I|+1}(u,\mathbf d_I)
 \mathcal U_{g_2,|J|+1}(v,\mathbf d_J).
\end{aligned}\label{eq:H9}
\end{equation} Here $\mathbf d_{\widehat{1,j}}$ denotes the tuple $\mathbf d$ with the entries $d_1$ and $d_j$ removed, and $\mathbf d_{\widehat1}$ the tuple with $d_1$ removed. The indices $u$ and $v$ are the degrees at the two new markings corresponding to the branches of the node. The last split is ordered. We call the amplitudes on the right of \eqref{eq:H9} the \emph{children} of the amplitude on the left: a merging or nonseparating child has complexity $h-1$, and the two factors of a separating term have complexities summing to $h-1$. Its initial values, computed directly from
\eqref{eq:H6}, are \begin{equation}
\mathcal U_{0,3}(0,0,0)=1,\qquad
\mathcal U_{1,1}(1)=\frac18,\qquad
\mathcal U_{1,1}(0)=\frac{\theta_1}{8}.
\label{eq:H10}
\end{equation} Equations \eqref{eq:H9}--\eqref{eq:H10}, with negative degrees and unstable factors set
to zero, determine all amplitudes by increasing \(h\).

By \eqref{eq:P11}, \begin{equation}
1=\theta_0\le\theta_1\le\cdots\le C_\theta
\label{eq:H11}
\end{equation} uniformly in \(Y\in[0,y_0]\). Thus every term of \eqref{eq:H9} is nonnegative,
and the continuous endpoint \eqref{eq:H3} defines a positive limiting theory.

\subsection{The harmonic model}\label{harmonic-model-and-the-needed-uniform-coefficient-estimates}

We now introduce the explicit model with which the amplitudes will be
compared. It depends on the degrees only through the deficit \(m\), and it
is chosen so that the linear part of~\eqref{eq:H9} reproduces it exactly
on average; this is the mass identity~\eqref{eq:H13} below, and the
reason for the name: the model is harmonic for the linear part of the
recursion. At the
endpoint \(Y=y_0\) the model is the denominator in the normalization
\eqref{eq:H4} of pure \(\psi\) intersection numbers (see~\eqref{eq:PF13}),
so the comparison theorem extends the coarse bounds on pure \(\psi\)
intersections to the whole parameter interval.

Suppress \(Y\) temporarily and put \begin{equation}
A_m(h)=[z^m]\bar H(z)^{-h},\qquad
\mathcal M_{g,n}(\mathbf d)=3^h\Gamma(h)4^{-g}A_m(h),\qquad
R=\mathcal U/\mathcal M.
\label{eq:H12}
\end{equation} All \(A_m(h)\) are positive. Define \[
q_\ell=\theta_\ell\,\frac{A_{m-\ell}(h-1)}{A_m(h)},
\qquad S_0=\sum_{\ell=0}^m q_\ell,\qquad
L=2D+n=3h-2m.
\] For \(h\ge2\), the identity \(\bar K=\bar H+(2/3)z\bar H'\) gives exactly \begin{equation}
\sum_{\ell=0}^m q_\ell
\left(1-\frac{2(m-\ell)}{3(h-1)}\right)=1.
\label{eq:H13}
\end{equation} Indeed, the generating series of the left side, before division by
\(A_m(h)\), is \[
\bar K^{-1}\left(\bar H^{-(h-1)}
-\frac{2z}{3(h-1)}(\bar H^{-(h-1)})'\right)
=\bar K^{-1}\bar K\bar H^{-h}=\bar H^{-h}.
\]

Dividing \eqref{eq:H9} by the model turns its linear terms into a
weighted sum of the ratios \(R\) at the children. We check that these
weights have total mass one when the pivot degree \(d_1\) is positive;
this is the meaning of \eqref{eq:H13}. Assume \(g\ge2\); states of genus at
most one are terminal in the argument of Section~\ref{uniform-comparison}.
A merging child has type \((g,n-1)\) and a nonseparating child type
\((g-1,n+1)\); both have complexity \(h-1\) and deficit \(m-\ell\). Hence the
model ratio of a merging child to its parent is
\(A_{m-\ell}(h-1)/(3(h-1)A_m(h))\), and that of a nonseparating child is four
times larger, because of the factor \(4^{-g}\) in \eqref{eq:H12}. The
merging children carry the coefficients \((2d_j+1)\theta_\ell\), \(j\ge2\),
and there are \(d_1+\ell-1\) ordered nonseparating pairs, each with
coefficient \(\theta_\ell/2\). Since \(2D+n=3h-2m\), the total weight is
\[
 \sum_{\ell=0}^m q_\ell\,
 \frac{\sum_{j\ge2}(2d_j+1)+2(d_1+\ell-1)}{3(h-1)}
 =\sum_{\ell=0}^m q_\ell\,\frac{3(h-1)-2(m-\ell)}{3(h-1)}=1
\]
by \eqref{eq:H13}. Thus, for a positive pivot, the linear model weights form
a probability kernel \(P\), and positivity of the separating terms gives
\(R\ge PR\).

At an all-zero state, where every degree and hence the pivot vanishes,
the formal \(\ell=0\) count includes one
nonseparating pair with multiplicity \(-1\), and also includes \(n-1\)
invalid mergers. The true linear mass is instead \begin{equation}
s=1+\frac{3-n}{3(h-1)}
       \frac{A_m(h-1)}{A_m(h)}.
\label{eq:H14}
\end{equation} Since \(\bar H(0)=1\) and \(\bar H^{-1}\) has positive coefficients,
\(A_m(h-1)\le A_m(h)\). In the range \(g\ge2\), therefore, \begin{equation}
s\ge\frac23,\qquad
s\le1\ (n\ge3),\qquad s\le1+O(h^{-1})\ (n=1,2).
\label{eq:H15}
\end{equation} Thus for a zero pivot the linear terms need not have mass exactly one.

We now record the analytic information needed below. The function
\begin{equation}
P_Y(w)=\frac{\bar H_Y(1)}{\bar H_Y(w)}
\label{eq:H16}
\end{equation} is the probability generating function of
\eqref{eq:P12}, recalled here. By
Subsection~\ref{compact-normalization-and-uniform-lattice-estimates}, it has
mean \(3/2\) and, uniformly in the compact parameter interval, it is
analytic on a disk of radius \(1+\eta>1\); its constant and linear
coefficients are bounded below by positive constants; and its variance
is bounded above and away from zero. Let \(S_h\) be the sum of \(h\)
independent variables with probability generating function
\eqref{eq:H16}. The uniform local limit theorem \eqref{eq:P13}, proved
there by Fourier inversion, gives \begin{equation}
\Pr(S_h=m)\ge c_Bh^{-1/2}
\quad\text{when }|m-3h/2|\le B\sqrt h,
\label{eq:H17}
\end{equation} and at bounded displacement, where the odd cubic term
integrates to zero, \begin{equation}
\Pr(S_h=3h/2+d)
=\frac{\mathrm{e}^{-d^2/(2h\sigma_Y^2)}}{\sqrt{2\pi h\sigma_Y^2}}
  (1+O_B(h^{-1})),\qquad |d|\le B,
\label{eq:H18}
\end{equation} at integer arguments. The same Fourier argument applies after tilting at a
positive radius \(\rho\) in a compact subset of the analytic interval; here $\rho$ is a radius in the normalized $z$-coordinate. If
\(\mu_Y(\rho)=-\rho\bar H_Y'(\rho)/\bar H_Y(\rho)=m/h\), it gives \begin{equation}
A_m(h)=
\frac{\bar H_Y(\rho)^{-h}\rho^{-m}}{\sqrt{2\pi h\,\rho\mu_Y'(\rho)}}
 (1+O(h^{-1}))
\label{eq:H19}
\end{equation} uniformly on such compact tilted families. Near the mean, the tilted
rate function has bounded second derivative and vanishing value and
first derivative; hence, for a fixed sufficiently small \(c_0>0\), \begin{equation}
\Pr(S_h=3h/2-u)\ge ch^{-1/2}\mathrm{e}^{-Cu^2/h},
\qquad 0\le u\le c_0h.
\label{eq:H20}
\end{equation}

The coefficient/probability relation is \begin{equation}
A_m(h)=\bar H(1)^{-h}\Pr(S_h=m).
\label{eq:H21}
\end{equation} Consequently, when \(L\le\sqrt h\), \eqref{eq:H11}, \eqref{eq:H17}, and the sum of the
numerator probabilities give \(S_0\le C\sqrt h\). When \(L\ge6\), the
minimum factor in \eqref{eq:H13} is \((L-3)/(3(h-1))\), so \(S_0\le Ch/L\).
Combining the two estimates, \begin{equation}
S_0\le C\sqrt h,\qquad S_0\le Ch/L\quad(L\ge6).
\label{eq:H22}
\end{equation} All constants here are uniform in \(Y\), and the finitely many states
of bounded complexity are absorbed into them.
